\section{Numerical Experiments}
\label{sec:experiments}

The experiments will distinguish theorem verification, bound usefulness, dimensional reduction, and downstream solver behavior.

\subsection{Exhaustive small-instance verification}

Small systems will use \(L\in\{2,3,4\}\) and small \(n_l\). For every instance, the study will enumerate all FIFO-feasible sequences and all contiguous platoon partitions to compute
\begin{equation}
D^*,\qquad D^*_{\Pi},\qquad G(\Pi),\qquad B(\Pi).
\end{equation}
The required validity check is
\begin{equation}
0\le G(\Pi)\le B(\Pi).
\end{equation}
The tightness measures are
\begin{align}
\text{Tightness ratio}&=\frac{G(\Pi)}{B(\Pi)},\\
\text{Bound excess}&=B(\Pi)-G(\Pi).
\end{align}
Cases with \(B(\Pi)=0\) will be reported separately.

\subsection{Dimension-reduction comparisons}

The following methods will be compared:
\begin{enumerate}[label=(\roman*)]
  \item vehicle-level scheduling without platooning;
  \item critical-headway platooning;
  \item fixed-threshold and fixed-size platooning; and
  \item the proposed bound-aware partition.
\end{enumerate}
Comparisons will control either the theoretical loss budget or the downstream model dimension.

\subsection{Scalability and traffic conditions}

The large-instance experiments will vary:
\begin{itemize}
  \item the number of approaches and vehicles;
  \item balanced and unbalanced approach demand;
  \item Poisson, uniform, and bursty arrivals;
  \item the ratio \(\hS/\hF\);
  \item the maximum platoon size; and
  \item downstream solver time limits.
\end{itemize}

Reported computational measures will include the number of platoons, number of ordering variables, Gurobi runtime, branch-and-bound node count, time to first feasible solution, and objective quality at fixed time limits.

\subsection{Pending implementation decisions}

\begin{itemize}
  \item Exact MILP representation of the partition decision.
  \item Discrete sets of acceptable loss budgets \(\varepsilon\) and size budgets \(\overline C\).
  \item Number of Monte Carlo replications and confidence intervals.
  \item Whether the original vehicle-level model can prove optimality within each time limit.
\end{itemize}
